% Keep the overview as a map; focus on the batch passed to distillation.
\begin{frame}{Challenges}
  \centering
  \setbeamercolor{itemize item}{fg=CVQCTeal}
  \begin{tikzpicture}[>={Latex[length=2mm,width=1.4mm]}]
    \path[use as bounding box] (-6.6,-3.20) rectangle (6.6,3.05);

    \begin{scope}[shift={(-3.2,0)}]
      \path[fill=CVQCPale,rounded corners=4mm]
        (-3.4,-3.15) rectangle (3.4,2.15);
      \node[anchor=north west,inner sep=0pt,
        font=\small\bfseries,text=CVQCInk] at (-3.10,1.95)
        {$\mathsf{QPIP}_1$};
      \protocolflowchart[fill=CVQCGoldPale,rounded corners=1.5mm,
        inner xsep=5pt]
    \end{scope}

    \node[anchor=west,inner sep=0pt,text=CVQCInk] at (0.80,0.25) {%
      \begin{minipage}{5.7cm}
        \raggedright
        \begin{itemize}
          \setlength{\itemsep}{1.0em}
          \item What does MBQC guarantee?
          \item What is distillable?
        \end{itemize}
      \end{minipage}%
    };
    \draw[->,very thick,draw=CVQCTeal,shorten >=2pt]
      (0.50,0.25) -- (rawmagic.east);
  \end{tikzpicture}

  \note{
    \footnotesize
    The highlighted batch connects preparation to distillation. Ask what
    MBQC and its trap-copy tests guarantee against a malicious prover,
    and whether that joint guarantee suffices for distillation.
    Recall the earlier example verbally. With probability
    $p=1/\mathrm{poly}$, all copies may be maliciously bad together;
    otherwise all are ideal. Each marginal has error at most p, but this
    alone does not establish distillability. If that branch passes with
    constant probability and gives a wrong output, its accepted bad weight
    remains inverse-polynomial. We are not claiming every all-bad batch
    passes: reliable rejection would eliminate the branch.
    We instead need the required factorizable error structure and a joint
    bound on how many copies are bad, except for negligible accepted weight.
    Authentication addresses malicious execution; it does not certify
    correctly prepared magic states. Next, explain the MBQC preparation
    and its error structure. Keep the counterexample off the slide.
  }
\end{frame}
